\chapter{能力二}
\section{无偏估计}
很好的问题，这是入门时最容易懵的地方。

\textbf{估计量的定义就是：样本的任意函数。}

你有数据 $X_1, X_2, \ldots, X_n$，你想猜 $\mu$。只要是用这些数据算出来的一个数，就叫估计量，没有任何其他限制。

比如 $n=3$，数据是 $2, 4, 6$：

\begin{center}
\begin{tabular}{lll}
\hline
估计量 & 公式 & 算出来 \\
\hline
$\bar{X}$ & $(2+4+6)/3$ & $4$ \\
$X_1$ & 只取第一个 & $2$ \\
$(X_1+X_3)/2$ & $(2+6)/2$ & $4$ \\
常数 $5$ & 不管数据，直接报 $5$ & $5$ \\
$\max(X_i)$ & 取最大值 & $6$ \\
\hline
\end{tabular}
\end{center}

这五个都是"合法"的估计量——你都可以拿来交给别人说"我估计 $\mu = $ 这个数"。

\textbf{问题就在这里}：到底用哪个？

直觉上 $\bar{X}$ 最好，但"直觉"不是数学。能力二要做的事就是：给出有数学保证的标准，说清楚为什么 $\bar{X}$ 比 $X_1$ 好，好在哪里，好多少。

一句话总结：\textbf{估计量 $=$ 数据的函数，函数有无穷多个，所以估计量不唯一。}能力二就是在这无穷多个里面找"最好的"。
\begin{note}
    总结一句话：你只做一次实验，\textbf{估计量}是你选的公式，\textbf{估计值}是代入数据后算出的那个数。``多次实验''是数学上验证公式好不好用的思想实验，不是真的要做很多次。
\end{note}
\begin{note}
    \textbf{第一个：不知道 $\mu$ 怎么对准？}

``对准''不是你去瞄准一个已知的靶，而是通过代数推导证明一个必然成立的性质。

$X_1, X_2, \ldots, X_n$ 来自 $N(\mu, \sigma^2)$，所以 $E[X_i] = \mu$，于是：

$$E[\bar{X}] = \frac{1}{n}\sum_{i=1}^n E[X_i] = \frac{1}{n} \cdot n\mu = \mu$$

这里 $\mu$ 始终是个符号，从没有变成具体数字。这个推导对 $\mu = 3$ 成立，对 $\mu = 100$ 也成立，对任何 $\mu$ 都成立。

就好像你证明了``一把没有系统误差的尺子，量出来的平均值等于真实长度''——不需要知道真实长度是多少，这是尺子本身的结构性质。
\end{note}
x1，x2都是一次实验里面的个体
\begin{dxtips}
也就是说因为他服从这个分布所以不管选第一次实验x1，第二次实验x2都要把他的均值当作真实正态分布的均值对吗。而不是选取一次实验的最大值，最小值直接当他服从的正态分布的均值。而是选取他的平均值充当
\end{dxtips}
\begin{note}[UMVUE 解决的是选公式的问题]
你确实只做一次实验，得到一个数。但在做实验之前，你要选公式。

方差小的公式，代入数据后得到的数大概率更靠近真值；方差大的公式，这一次的结果可能差很远。

类比：你只射一箭，但要事先选弓。弓A精度高，弓B精度差。虽然只射一次，但选哪把弓决定了这一箭大概率落在哪里。UMVUE 就是在所有无偏的弓里，选精度最高的那把。

所以 UMVUE 是\textbf{选公式阶段}的理论，不是做完实验之后的事情。
\end{note}

\begin{example}[UMVUE 的具体比较]
设 $X_1, X_2, X_3, X_4 \overset{\text{i.i.d.}}{\sim} N(\mu, 1)$，估计 $\mu$。三个候选公式均无偏（期望均等于 $\mu$）：

$$\hat{\mu}_A = \bar{X} = \frac{X_1+X_2+X_3+X_4}{4}, \quad \mathrm{Var}(\hat{\mu}_A) = \frac{1}{4}$$

$$\hat{\mu}_B = \frac{X_1+X_2}{2}, \quad \mathrm{Var}(\hat{\mu}_B) = \frac{1}{2}$$

$$\hat{\mu}_C = X_1, \quad \mathrm{Var}(\hat{\mu}_C) = 1$$

方差衡量的是：这个公式在不同数据下的结果散布多广。$\bar{X}$ 的结果永远集中在 $\mu$ 附近，$X_1$ 的结果可能飘得很远。

在所有对 $\mu$ 的无偏估计量中，$\bar{X}$ 方差最小，故 $\bar{X}$ 是 $\mu$ 的 UMVUE。
\end{example}
\begin{note}
    这是统计里一个让人头疼的惯例，其实它们是两个不同的东西：

\begin{itemize}
  \item \textbf{实验前：$X_1$ 是随机变量（大写）}
  
  还不知道会观测到什么值，可能是 $160$，可能是 $180$，服从 $N(\mu, \sigma^2)$。这时候可以谈期望：$E[X_1] = \mu$。

  \item \textbf{实验后：$x_1$ 是观测值（小写）}
  
  已经量完了，是一个固定的数，比如 $172$。这时候不谈期望，直接代入算：$\bar{x} = (172 + 168 + \cdots)/n$。
\end{itemize}

因为教材里经常两个都写大写 $X_1$，没有严格区分。区分的方法就是看语境：出现 $E[\cdot]$ 就当随机变量（实验前）；出现具体数字代入就当观测值（实验后）。


\end{note}
\begin{dxtips}
    找无偏估计量永远发生在实验前，与具体采样永远无关。需要把所有可能出现的样本进行 $h(X)$ 的映射，再与其出现的概率做乘积，求和得到期望：
$$E[h(X)] = \sum_k h(k) \cdot P(X=k) = g(\theta)$$
这个期望等于我们想估计的目标 $g(\theta)$，即为无偏。
\end{dxtips}
\begin{note}
大写 $X_1$ 是``还没揭晓的牌''，小写 $x_1$ 是``已经翻开的牌''。期望只对没揭晓的牌讲。
\end{note}
\begin{dxtips}
    

所以无偏估计量是找一个公式，用这个公式代入所有可能的样本后，按照这些样本的分布分配给它们概率，然后这个期望等于我们想估计的东西。
\end{dxtips}
\begin{dxtips}
    证明没有无偏函数可以从解析性和多项式来看
\end{dxtips}
\subsection{顺序统计量}
```latex
\begin{example}{为什么需要顺序统计量：从电路寿命说起}
你有 $n = 5$ 个灯泡，串联组成一条电路——\textbf{一个坏了，全线中断}。
问：这条电路能撑多久？

你知道每个灯泡的寿命 $X_i$ 服从同一分布 $F(x)$，也可以算出样本均值 $\bar{X}$。
但 $\bar{X}$ 在这里毫无用处：它告诉你的是"平均每个灯泡能用多久"，
而电路的寿命取决于\textbf{最先坏掉的那个}，即
\[
    T_{\text{串联}} = \min\{X_1, X_2, X_3, X_4, X_5\} = X_{(1)}.
\]
如果改成并联电路——\textbf{全部坏了才断}，则寿命取决于\textbf{最后一个}：
\[
    T_{\text{并联}} = \max\{X_1, X_2, X_3, X_4, X_5\} = X_{(n)}.
\]

\begin{itemize}
    \item $\bar{X}$ 描述数据的"整体水平"，把每个点的位置信息\textbf{平均掉了}。
    \item $X_{(1)}, \dots, X_{(n)}$ 保留了"第几名"的位置信息，
          专门回答极值、中间值、排名类问题。
\end{itemize}

凡是问题出现"最大""最小""第 $r$ 小""中位数"，本质上都是在问顺序统计量——
这就是它被发明的根本原因。
\end{example}
\begin{dxtips}
    均值方差是"融合"数据，顺序统计量是"保留位置"的数据。
    一类问题问平均水平用前者；一类问题问第几名用后者。
    两种工具解决两类完全不同的问题。
\end{dxtips}
\begin{note}[顺序统计量]
顺序统计量在一个特定情况下才出现：\textbf{参数出现在分布的支撑边界上（即参数决定数据的范围）。}

\textbf{用顺序统计量：}
\begin{itemize}
  \item $U(\theta_1, \theta_2)$：数据必须在 $[\theta_1, \theta_2]$ 里，参数就是边界 $\rightarrow$ 用 $X_{(1)}$、$X_{(n)}$
  \item $U(0, \theta)$：数据在 $[0, \theta]$ 里，$\theta$ 是上界 $\rightarrow$ 用 $X_{(n)}$
  \item 双参数指数分布 $f(x) = e^{-(x-\theta)}$，$x \geq \theta$：$\theta$ 是数据的下界 $\rightarrow$ 用 $X_{(1)}$
\end{itemize}

\textbf{不用顺序统计量：}
\begin{itemize}
  \item $N(\mu, \sigma^2)$：数据范围是整个实数轴，$\mu$、$\sigma^2$ 出现在密度公式里 $\rightarrow$ 用 $\bar{X}$、$s^2$
  \item $P(\lambda)$：$\lambda$ 出现在概率公式里，不影响范围 $\rightarrow$ 用 $\bar{X}$
\end{itemize}
\begin{dxtips}
    设 $X_1, X_2, \cdots, X_8$ i.i.d. $\sim U(0, \theta)$，
观测值为 $1, 2, 3, 4, 5, 6, 7, 8$，
则 $X_{(8)} = \max\{X_1, \cdots, X_8\} = 8$，
用 $X_{(8)}$ 估计 $\theta$。
我们求 $X_{(n)}$ 的期望来估计 $\theta$，
就像我们拿 $\bar{X}$ 的期望估计正态分布的 $\mu$。
求期望里面都是变量不是真实的数字是实验前。
\end{dxtips}
\begin{example}
$X_{(n)}$ 虽然是从 $X_1,\dots,X_8$ 里来的，但它不是随机抽的一个，
它是特意挑出最大的那个。挑最大值这个动作改变了分布。

具体来说，设 $\theta = 10$：

$$P(X_i \leq 5) = \frac{5}{10} = 0.5$$

随便一个 $X_i$ 有 $50\%$ 概率 $\leq 5$，合理。

$$P(X_{(8)} \leq 5) = \left(\frac{5}{10}\right)^8 = \frac{1}{256} \approx 0.4\%$$

8 个数里最大值还 $\leq 5$，概率只有 $0.4\%$。

直觉：8 个数都取到 5 以下，这件事几乎不可能发生，
因为总会有某个数跑到 5 以上。
所以 $X_{(n)}$ 虽然也是 $X_i$ 里的一个，
但它天生偏大，分布和普通 $X_i$ 不一样。
\end{example}
\begin{dxtips}
    这里的 $X_1, \dots, X_8$ 不是最大值 $X_{(8)}$，而是八个原始样本。
最大值小于 $5$ 等于全部八个样本都小于 $5$ 的概率，
每一个都服从 $U(0,\theta)$，小于 $5$ 的概率都是 $\dfrac{5}{\theta}$。
\end{dxtips}
\begin{note}
$$F_{(n)}(x) = \left(\frac{x}{\theta}\right)^n$$

$$f_{(n)}(x) = \frac{nx^{n-1}}{\theta^n}$$

$$E[X_{(n)}] = \int_0^\theta x \cdot \frac{nx^{n-1}}{\theta^n}\,dx 
= \frac{n}{\theta^n}\int_0^\theta x^n\,dx 
= \frac{n}{\theta^n} \cdot \frac{\theta^{n+1}}{n+1} 
= \frac{n\theta}{n+1}$$
\end{note}

\begin{dxtips}
    原本从 $0$ 积分到 $\theta$，故
$$E[X_{(n)}] = \frac{n\theta + 0}{n+1}$$
现从 $\theta_1$ 积分到 $\theta_2$，把 $\theta$ 换成 $\theta_2$，$0$ 换成 $\theta_1$：
$$E[X_{(n)}] = \frac{n\theta_2 + \theta_1}{n+1}$$

同理，$E[X_{(1)}]$ 把积分区域反过来，从 $\theta_2$ 积分到 $\theta_1$，
$n$ 后面跟着的是 $\theta_1$，另外加的是起始点 $\theta_2$：
$$E[X_{(1)}] = \frac{n\theta_1 + \theta_2}{n+1}$$
\end{dxtips}
\begin{note}

参数在密度公式里 $\rightarrow$ 用 $\bar{X}$、$s^2$ 这类自然统计量；参数在数据的取值范围上 $\rightarrow$ 用极值 $X_{(1)}$、$X_{(n)}$。
\end{note}
\begin{note}[顺序统计量期望的等间距直觉]
$n$ 个点均匀撒在 $[\theta_1, \theta_2]$ 上，平均来说这 $n$ 个点会把区间切成 $n+1$ 段等长的小段，每段长度为：
$$\frac{\theta_2 - \theta_1}{n+1}$$

因此：
\begin{itemize}
  \item $X_{(1)}$ 平均落在距离 $\theta_1$ 一段的位置：
  $$E[X_{(1)}] = \theta_1 + \frac{\theta_2-\theta_1}{n+1} = \frac{n\theta_1+\theta_2}{n+1}$$

  \item $X_{(n)}$ 平均落在距离 $\theta_2$ 一段的位置：
  $$E[X_{(n)}] = \theta_2 - \frac{\theta_2-\theta_1}{n+1} = \frac{\theta_1+n\theta_2}{n+1}$$
\end{itemize}

$X_{(1)}$ 是最小值，靠近 $\theta_1$，故 $\theta_1$ 权重大（$n$ 份），$\theta_2$ 权重小（$1$ 份）；$X_{(n)}$ 反之。

记忆：$n$ 个点把区间切成 $n+1$ 段，最小值在第 $1$ 段末，最大值在第 $n$ 段末。
\end{note}
\end{note}
\section{umvue+CR}
\subsection{fisher信息}
\begin{note}
\textbf{Fisher 想解决的问题：}这组数据对 $\theta$ 有多"敏感"？

如果 $\theta$ 稍微变一点，数据出现的概率就大幅变化——说明数据能精确分辨
$\theta$ 的位置，信息量大。如果怎么改 $\theta$ 似然函数都几乎不动，
说明数据对 $\theta$ 几乎无感，信息量小。
\end{note}

\begin{note}
得分函数 $S = \dfrac{\partial}{\partial\theta}\ln f(X;\theta)$ 正好衡量这件事：
它就是对数似然关于 $\theta$ 的导数，即似然函数对 $\theta$ 变化的敏感程度。
\begin{itemize}
  \item $S$ 大 $\to$ 似然对 $\theta$ 变化剧烈 $\to$ 数据信息多
  \item $S$ 小 $\to$ 似然对 $\theta$ 几乎无感 $\to$ 数据信息少
\end{itemize}
\end{note}
\begin{note}
$f(X;\theta)$ 是什么：是单个 $X$ 的概率密度函数

就是概率密度函数，但把视角反过来——数据 $X$ 已经观测到了，$\theta$ 是未知的。
$$f(X;\theta) = \text{``如果参数是 } \theta \text{，观测到这个 } X \text{ 的概率有多大？''}$$
这就是似然函数，$\theta$ 变，这个概率跟着变。
\begin{itemize}
  \item 单个 $X$：得分函数为 $S = \dfrac{\partial}{\partial\theta}\ln f(X;\theta)$，
        只用 $f(X;\theta)$，与 $L(\theta)$ 无关。
  \item 多个 $X_1,\dots,X_n$：似然函数 $L(\theta)=\prod_{i=1}^n f(X_i;\theta)$，
        得分函数为
        $$S_n = \frac{\partial}{\partial\theta}\ln L(\theta)
        = \sum_{i=1}^n \frac{\partial}{\partial\theta}\ln f(X_i;\theta)
        = \sum_{i=1}^n S(X_i;\theta)$$
        此时得分函数与 $L(\theta)$ 产生联系。
\end{itemize}
\end{note}
\begin{dxtips}
    ln是为了累乘变累加好算，对θ求导数因为本来就是腰算他的变化率
\end{dxtips}
\begin{note}
    似然函数英文叫 Likelihood function。
\textbf{Likelihood} 是"可能性/合理程度"的意思，不是严格的概率，
而是"这个 $\theta$ 值有多合理？"
中文"似然"就是翻译自 likelihood——似乎如此，好像这个参数值是对的。
\end{note}
\begin{itemize}
  \item $f(x \mid \theta)$：竖线后面是条件，贝叶斯里 $\theta$ 被当作随机变量，
        "在 $\theta$ 已知的条件下"
  \item $f(x;\theta)$：分号前面是\textbf{观测值}（已知），
        分号后面是\textbf{参数}（待估），频率派里 $\theta$ 不是随机变量，只是未知常数
\end{itemize}
\begin{note}
得分函数 $S(X;\theta)$ 本身是随机的，因为 $X$ 是随机的。

对于固定的 $\theta$，不同的 $X$ 值会给出不同的 $S$：
$$S(X;\theta) = \frac{\partial}{\partial\theta}\ln f(X;\theta)$$
$X$ 变，$S$ 就变。所以 $S$ 不是一个数，是一个随机变量。

我们已经知道 $E[S]=0$，均值没有信息。因为 $f(x;\theta)$ 是关于 $x$ 的概率密度，
积分恒为 $1$，两边对 $\theta$ 求导所以等于 $0$。

那用什么来总结"$S$ 平均能偏离 $0$ 多远"？自然就是方差：
\begin{itemize}
  \item $\operatorname{Var}(S)$ 大 $\to$ $S$ 经常取大值 $\to$ 似然对 $\theta$ 平均很敏感 $\to$ 信息多
  \item $\operatorname{Var}(S)$ 小 $\to$ $S$ 总在 $0$ 附近 $\to$ 似然对 $\theta$ 几乎无感 $\to$ 信息少
\end{itemize}
\end{note}
\begin{note}
    $E[S]=0$ 的推导依赖一步交换：
$$\frac{\partial}{\partial\theta}\int f(x;\theta)\,dx = \int \frac{\partial}{\partial\theta}f(x;\theta)\,dx$$
左边对 $\int f\,dx = 1$ 求导得 $0$，右边整理后就是 $E[S]$，所以 $E[S]=0$。

但这个交换要求\textbf{积分上下限不含 $\theta$}。$U(0,\theta)$ 的积分上限就是 $\theta$，交换不合法：
$$\frac{\partial}{\partial\theta}\int_0^\theta \frac{1}{\theta}\,dx = \frac{\partial}{\partial\theta}1 = 0$$
$$\int_0^\theta \frac{\partial}{\partial\theta}\frac{1}{\theta}\,dx = \int_0^\theta -\frac{1}{\theta^2}\,dx = -\frac{1}{\theta}$$
两边不相等，所以 $E[S] = -\dfrac{1}{\theta} \neq 0$，C-R 下界对 $U(0,\theta)$ 直接失效。
\end{note}
\subsection{CR}
\begin{note}
\textbf{C-R 下界在说什么：}

对任意无偏估计量 $\hat{\theta}$，方差不可能低于某个值：
$$\operatorname{Var}(\hat{\theta}) \geq \frac{1}{nI(\theta)}$$
就是给方差划了一条\textbf{不可逾越的底线}。
\end{note}

\begin{note}
\textbf{为什么被发明：}

能力二的核心问题是"哪个估计量更好"，判断标准是方差小。
但比较 A 和 B，你只知道 A 比 B 好，不知道 A 是不是最好的。

C-R 下界解决了这个问题——给你一个\textbf{绝对基准}：
\begin{itemize}
  \item $\operatorname{Var}(\hat{\theta}) = \dfrac{1}{nI(\theta)}$ $\to$ 最优，没有比它更好的无偏估计量
  \item $\operatorname{Var}(\hat{\theta}) > \dfrac{1}{nI(\theta)}$ $\to$ 还\textbf{可能}有改进空间
\end{itemize}
\end{note}
\newpage
\begin{proof}
    \textbf{第一步：}无偏条件 $E[\hat{\theta}] = \theta$，两边对 $\theta$ 求导：
$$\int \hat{\theta}\cdot\frac{\partial f}{\partial\theta}\,dx = 1$$
代入 $\textcolor{red}{\dfrac{\partial f}{\partial\theta} = f\cdot S}$，得：
$$E[\hat{\theta}\cdot S] = 1$$

\textbf{第二步：}因为 $E[S]=0$，所以：
$$\operatorname{Cov}(\hat{\theta},\, S) = E[\hat{\theta}\cdot S] - E[\hat{\theta}]\cdot E[S] = 1$$

\textbf{第三步：}柯西-施瓦茨不等式：
$$\operatorname{Cov}(\hat{\theta},S)^2 \leq \operatorname{Var}(\hat{\theta})\cdot\operatorname{Var}(S) \implies 1 \leq \operatorname{Var}(\hat{\theta})\cdot I(\theta) \implies \operatorname{Var}(\hat{\theta}) \geq \frac{1}{I(\theta)}$$
$n$ 个样本再除以 $n$ 得 $\dfrac{1}{nI(\theta)}$。
\end{proof}
\begin{dxtips}
柯西-施瓦茨不等式才能把无偏估计的方差连接到不等式中。
我们已经知道无偏估计的期望是 $\theta$，为了出现 $\dfrac{\partial f}{\partial\theta}$
从而引入 $S$，所以对两边求导，再利用 $E[S]=0$ 的信息
和 $\operatorname{Var}(S) = I(\theta)$ 完成这个证明。
\end{dxtips}
\section{umvue}
\begin{note}
    \textbf{全称}：Uniformly Minimum Variance Unbiased Estimator，一致最小方差无偏估计量。

\textbf{为什么发明它}，从需求倒推：

你想估计 $\theta$，无偏是基本要求（平均不偏）。但满足无偏的估计量有无数个——$X_1$ 是无偏的，$\overline{X}$ 也是，随便凑个线性组合还是。\textbf{怎么选最好的那个？}

自然的标准：方差最小——离真值的随机误差最小。

但"最小"要小心：不能只对某个特定 $\theta$ 值最小，要对\textbf{所有} $\theta \in \Theta$ 同时最小，这就是"Uniformly"（一致）的含义。

所以 UMVUE 回答的问题是：
\begin{quote}
在全部无偏估计量里，有没有一个，不管真实参数取什么值，它的方差都比其他所有无偏估计量小？
\end{quote}
有的话，它就是 UMVUE，是无偏类里的理论最优。
\end{note}
\begin{dxtips}
    uniformity 一致的，一致收敛是uniform，m就是最小 v就是方差 ue就是无偏估计
\end{dxtips}
\subsection{完备充分统计量}
\begin{note}[充分统计量不能丢信息]
    用正态 $N(\mu, 1)$，$n=2$ 来说，取\textbf{不充分}的统计量 $T = X_1$。

你观测到 $X_1 = 3,\ X_2 = 5$。

\textbf{用 $T = X_1 = 3$ 去估计 $\mu$}：你只看第一个观测值，猜 $\mu \approx 3$。
但 $X_2 = 5$ 还摆在那里，它也服从 $N(\mu, 1)$，里面明显还有 $\mu$ 的信息——
看了 $X_2 = 5$ 之后你会更新判断，觉得 $\mu$ 应该在 4 附近。
$T = X_1$ 把 $X_2$ 里的 $\mu$ 信息\textbf{扔掉了}，这就是"丢了 $\theta$ 信息"。

\textbf{用 $T = \overline{X} = 4$ 就不丢}：
知道了 $\overline{X} = 4$，再去看 $X_1 = 3,\ X_2 = 5$ 的具体分法——
$X_1$ 比均值低 1，$X_2$ 比均值高 1——这件事和 $\mu$ 是多少完全没关系，只是随机波动。
没有任何新的 $\mu$ 信息可以被挖出来了。

也就是说，知道了 $\overline{X} = 4$，不用再知道到底是 $(3,5)$ 还是 $(4,4)$，
因为对"$\mu$ 大概是多少"这个问题，两种情况提供的信息量完全一样。

\textbf{一句话}：$T$ 是不是充分的，就看"知道 $T$ 之后，样本剩下的部分还能不能更新你对 $\theta$ 的判断"。能更新就说明丢了信息，不充分；不能更新就是充分的。
\end{note}
\begin{note}
\textbf{完备空间类比}

完备空间的核心是：极限不会"跑出去"——每个 Cauchy 列的极限都在空间里，没有"洞"。

$\mathbb{Q}$ 不完备：$1, 1.4, 1.41, 1.414, \ldots$ 是 Cauchy 列，极限 $\sqrt{2}$ 跑到 $\mathbb{Q}$ 外面去了。

$\mathbb{R}$ 完备：任何 Cauchy 列的极限都还在 $\mathbb{R}$ 里，无处可逃。

\textbf{完备充分统计量的逻辑完全一样：}

不完备的 $T$：存在非零函数 $g(T)$，它的期望"跑到零去了"（$E[g(T)]=0\ \forall\theta$），但 $g$ 本身不是零。
就像 $\sqrt{2}$ 逃出了 $\mathbb{Q}$——非零的东西期望上"假装是零"。

完备的 $T$：期望等于零这件事无处可逃，只要 $E[g(T)]=0$，$g$ 就必须是零本身，没有任何非零函数能"藏"在零期望后面。

\textbf{为什么这样定义}

设计这个定义的动机就是逼出唯一性。假设 $g_1(T)$ 和 $g_2(T)$ 都是 $\theta$ 的无偏估计，那么：
$$E_\theta[g_1(T) - g_2(T)] = 0 \quad \forall\theta$$
完备性直接给出 $g_1(T) = g_2(T)$——基于 $T$ 的无偏估计\textbf{唯一}，唯一的那个自然是 UMVUE。定义就是为了这一步设计的。
\end{note}
\begin{dxtips}
    所以完备统计量就抱着基于T的无偏估计只有一个，因为只有一个所以一定是umvue对吗
\end{dxtips}